\lecture{Лекция 3}{}

\subsection{Методы одномерной минимизации: постановка}

Рассматривается задача
\[
  f(x) \longrightarrow \min, \qquad x \in [a,b] ,
\]
где $[a,b]$ — \textbf{исходный интервал неопределённости}.

\medskip
\noindent\textbf{Задача.} Локализовать точку минимума $x_*$ в как можно более
узком \textbf{конечном интервале неопределённости} $[a_k, b_k]$.

\subsubsection*{Откуда берётся одномерная задача}

В многомерной минимизации $\varphi(x) \to \min$, $x \in \R^n$, работает
двухшаговая схема: из начального приближения $x_0$ строится
последовательность
\[
  x_{k+1} = x_k + \alpha_k p_k , \qquad \varphi(x_{k+1}) < \varphi(x_k) ,
\]
причём на каждой итерации
\begin{enumerate}[label=\arabic*), leftmargin=2em]
  \item находим направление $p_k$;
  \item находим шаг $\alpha_k$.
\end{enumerate}

\inputfig{figures/fig-descent-step}

Поиск $\alpha_k$ — подзадача одномерной минимизации: минимизируется функция
одного переменного $\alpha \mapsto \varphi(x_k + \alpha p_k)$. Особый интерес
представляют методы, \textbf{не требующие подсчёта производной}: именно их мы и
разбираем.

\subsubsection*{Унимодальность}

\begin{definition}
Функция $f(x)$ \textbf{унимодальна} на $[a,b]$, если существует единственная
точка глобального минимума $x_*$ на $[a,b]$, причём
\[
  f \ \text{строго убывает на} \ [a, x_*) ,
  \qquad
  f \ \text{строго возрастает на} \ (x_*, b] .
\]
\end{definition}

\begin{idea}
Унимодальность — это ровно то свойство, которое позволяет по двум вычислениям
функции выбросить часть отрезка. Непрерывность или гладкость здесь не нужны:
всё держится на монотонности по обе стороны от минимума.
\end{idea}

\begin{lemma}\label{lem:unimodal}
Пусть $f$ унимодальна на $[a,b]$ и $x_1, x_2 \in [a,b]$, $x_1 < x_2$. Тогда
\[
  1^\circ\quad f(x_1) \ge f(x_2) \ \Longrightarrow\ x_* \in [x_1, b] ,
\]
\[
  2^\circ\quad f(x_1) < f(x_2) \ \Longrightarrow\ x_* \in [a, x_2] .
\]
\end{lemma}

\inputfig{figures/fig-unimodal}

\begin{proof}
$1^\circ$. Пусть, напротив, $x_* < x_1$. Тогда обе точки $x_1 < x_2$ лежат в
промежутке $(x_*, b]$, где $f$ строго возрастает, откуда $f(x_1) < f(x_2)$ —
противоречие с условием $f(x_1) \ge f(x_2)$. Значит $x_* \ge x_1$.

\medskip
$2^\circ$. Пусть, напротив, $x_* > x_2$. Тогда обе точки лежат в промежутке
$[a, x_*)$, где $f$ строго убывает, откуда $f(x_1) > f(x_2)$ — противоречие с
условием $f(x_1) < f(x_2)$. Значит $x_* \le x_2$.
\end{proof}

\subsubsection*{Как сравнивать методы}

\begin{idea}[title={Критерий сравнения}]
Методы одномерной минимизации сравниваются по \textbf{числу обращений к
вычислению функции}, требуемому для достижения заданной точности.

\medskip
Логика такая: сама $f$ обычно считается дорого (в многомерной схеме каждое
обращение — это вычисление $\varphi$ в новой точке $\R^n$), а вся арифметика
метода по сравнению с этим бесплатна. Поэтому метод тем лучше, чем сильнее он
сокращает интервал неопределённости \emph{на одно вычисление функции}.
\end{idea}
